\section{本章小结}
图像质量由入光、曝光、传感器读出和后续处理共同决定。曝光时间影响运动模糊，读出方式影响卷帘形变，增益与 Gamma 不能恢复已经丢失的信息。标定要串联“固定成像模式、多姿态采集、求解内参与畸变、独立验证、驱动加载、校正输出”。应区分原图与校正图的模型、单目标定板位姿与安装外参，并核对 ROS2 的 CameraInfo。镜头景深公式和专用 HDR 参数可以按需查阅。

\begin{problemset}
\item 基础：焦距 16 mm、入瞳直径 4 mm，光圈数为多少？
\item 基础：目标速度 600 pixel/s、曝光 5 ms，曝光内位移约多少？
\item 基础：全局快门是否能消除曝光过长造成的模糊？增益能否增加收到的光子数？
\item 基础：不考虑填充，100×100 的 Bayer8 与 RGB8 各占多少字节？
\item 基础：Gamma 改亮暗部后，能否保证恢复过曝区域的纹理？
\item 提高：标定时把 25 mm 方格误填为 25 m，会怎样影响内参和估计平移的单位尺度？
\item 提高：多相机曝光起点差 6 ms，目标速度 2 m/s，对应位置差多少？
\item 实践：按本章公式，用 Python 计算图像速度 1200 pixel/s、容许位移 0.6 pixel 时的最长曝光时间。
\item 基础：ROS2 标定界面的 CALIBRATE、SAVE、COMMIT 分别做什么？重启驱动后怎样确认标定结果仍生效？CameraInfo 中的 R 能否直接作为底盘外参？
\item 提高：训练照片的 RMS 很低，但校正后的边缘直线仍弯曲，应检查哪些数据与模型问题？对校正图做 PnP 时应使用哪组内参和畸变？
\end{problemset}

\section{本章验收}
\begin{enumerate}
    \item 在固定光照下拍摄同一运动目标，分别修改曝光和增益，提交图片及配置，解释清晰度与亮度的变化。
    \item 完成一次棋盘格标定，展示有效照片数、内参、畸变参数、重投影误差及一张未参与标定的验证照片。
    \item 用 ROS2 camera\_calibration 完成采集和保存，提交原始样本与 YAML；重启驱动后现场核对 CameraInfo，并用 image\_proc 展示原图和校正图。没有写入服务时演示手动文件加载。
    \item 结合传感器读出方式说明卷帘问题，列出至少两种噪声来源及相应排查方法。
\end{enumerate}
